% tikz-tensors-coords.tex -- (temporary) the commands that place by coordinates

% \tnchain[<pitch>]{<prefix>}{<from>}{<to>}{<style/label/side, ...>}
% Draw one line and put the tensors on it. Without <pitch> they are spread
% evenly from <from> to <to>, first and last on the ends; with it they start at
% <from> and step by <pitch>, and <to> only gives the direction. Site i becomes
% the node <prefix>i and the coordinate <prefix>i-leg, the point where its
% physical index leaves it and the point everything else connects to.
%
%   \tnchain{S}{(0,0)}{(4,0)}{tn triangle right/$A$/l, tn diamond/$C$/c,
%                             tn triangle left/$B$/r}
%
% Two coordinates for a row, and none per site: a row that is respaced, or
% turned, or given another tensor, is one line edited in one place.
%
% <side> is where the index leaves: `l' and `r' are the bottom corner of a
% triangle's flat side, pointing right and pointing left (corner 3 and corner 2 -- the two are one shape
% at a 180 degree turn, so no single index names both, and `south west' is on
% the edge above the vertex rather than on it), and `c' is the node's own south,
% which is what a box, a circle or a diamond takes.
% An anchor lands on the line that is drawn only because `tn node' sets
% outer sep=0pt. With TikZ's default, outer sep is half a line width and it is
% added to every anchor as well as to the border, so a corner anchor sits that
% far outside the stroke and an index placed on it runs beside the tensor's own
% edge instead of along it. Measured on a triangle pointing right: corner 3 at -7.86884pt
% with the default and -7.46884pt without, exactly half a width apart.
\def\tn@side@l{corner 3}
\def\tn@side@r{corner 2}
\def\tn@side@c{south}
\newcount\tn@cnt
\newcommand{\tnchain}[5][]{%
  \global\tn@cnt=0
  \foreach \tn@z in {#5}{\global\advance\tn@cnt by 1}%
  \foreach \tn@sty/\tn@lab/\tn@sd [count=\tn@i from 1] in {#5}{%
    \ifx\relax#1\relax
      \ifnum\tn@cnt>1
        \pgfmathsetmacro{\tn@t}{(\tn@i-1)/(\the\tn@cnt-1)}%
      \else \def\tn@t{0}\fi
      \coordinate (tn@spot) at ($#3!\tn@t!#4$);
    \else
      \pgfmathsetlength{\pgf@xa}{(\tn@i-1)*(#1)}%
      \coordinate (tn@spot) at ($#3!\the\pgf@xa!#4$);
    \fi
    \node[\tn@sty] (#2\tn@i) at (tn@spot) {\tn@lab};
    \coordinate (#2\tn@i-leg) at (#2\tn@i.\csname tn@side@\tn@sd\endcsname);
  }%
}

% \tnrow[<gap>]{<prefix>}{<origin>}{<style/label/side, ...>}
% Place tensors left to right, each <gap> clear of the one before it -- measured
% between their borders, not between their centres, so the spacing is right
% whatever shapes and sizes are in the row and stays right when one of them
% changes. \tnchain spreads a row evenly along a line you give it, which is what
% a chain of like tensors wants; this is for a row of unlike ones, where the
% distance that matters is the gap and working the centres out by hand is both
% tedious and the first thing to go stale.
%
%   \tnrow{P}{(0,0)}{tn box/$A$/c, tn rounded/$W$/c, tn diamond/$C$/c}
%
% Site i becomes the node <prefix>i and the coordinate <prefix>i-leg, exactly as
% in \tnchain.
\newdimen\tn@cur \newdimen\tn@rowy
\newcommand{\tnrow}[4][6mm]{%
  \coordinate (tn@rowstart) at #3;
  \pgf@process{\pgfpointanchor{tn@rowstart}{center}}%
  \global\tn@cur=\pgf@x \global\tn@rowy=\pgf@y
  \foreach \tn@sty/\tn@lab/\tn@sd [count=\tn@i from 1] in {#4}{%
    \node[\tn@sty, anchor=west] (#2\tn@i) at (\the\tn@cur,\the\tn@rowy) {\tn@lab};
    \coordinate (#2\tn@i-leg) at (#2\tn@i.\csname tn@side@\tn@sd\endcsname);
    \pgfextractx{\tn@x}{\pgfpointanchor{#2\tn@i}{east}}%
    \global\tn@cur=\tn@x \global\advance\tn@cur by #1\relax
  }%
}

% \tnlegs[<kind>]{<drop>}{<points>}
% A physical index from each point, every one ending <drop> below the lowest of
% them, so a row of legs is level whatever each started from -- a diamond hangs
% lower than a triangle's corner, and the drop is a rule rather than a y to be
% written out. <drop> is a length (16mm, 1.6cm), like \tngate's padding and
% unlike the line \tnchain is given, which is a pair of coordinates. Each leg's far end becomes <point>-tip, which is where a label
% for that index goes.
\newdimen\tn@x \newdimen\tn@y \newdimen\tn@lo \newdimen\tn@hi
\newcommand{\tn@spanx}[1]{% min and max x of a list of points
  \global\tn@lo=\maxdimen \global\tn@hi=-\maxdimen
  \foreach \tn@s in {#1}{%
    \pgfextractx{\tn@x}{\pgfpointanchor{\tn@s}{center}}%
    \ifdim\tn@x<\tn@lo \global\tn@lo=\tn@x \fi
    \ifdim\tn@x>\tn@hi \global\tn@hi=\tn@x \fi}}
\newcommand{\tn@lowy}[1]{% the lowest y of a list of points
  \global\tn@lo=\maxdimen
  \foreach \tn@s in {#1}{%
    \pgfextracty{\tn@y}{\pgfpointanchor{\tn@s}{center}}%
    \ifdim\tn@y<\tn@lo \global\tn@lo=\tn@y \fi}}
\newcommand{\tnlegs}[3][tn edge]{%
  \tn@lowy{#3}%
  \pgfmathsetlength{\pgf@yb}{\tn@lo-(#2)}\edef\tn@bot{\the\pgf@yb}%
  \begin{pgfonlayer}{tnedges}%
    \foreach \tn@n in {#3}{%
      \coordinate (\tn@n-tip) at (\tn@n |- 0,\tn@bot);
      % Start inside the tensor and let it cover the overlap -- nodes are drawn
      % over the edge layer. Begun exactly on the point, the leg's own stroke
      % ends half a line width short of where the outline's stroke ends and the
      % join reads as a gap; begun above it, the two merge. Upward only, so a
      % leg taken from the bottom corner of a triangle's flat side stays on that
      % side's own line and simply continues it down.
      \draw[#1] ($(\tn@n)+(0,2*\pgfkeysvalueof{/tn/line width})$) -- (\tn@n-tip);%
    }%
  \end{pgfonlayer}%
}

% \tngate[<padding>]{<node name>}{<label>}{<points>}{<drop>}
% A gate on those indices, <drop> below the lowest of them (a length, as in
% \tnlegs), as wide as they span plus <padding> at each end and centred on them. One point gives a small
% box, several a bar. It needs no legs of its own: a node is drawn over an
% edge, so a gate placed on indices that run past it interrupts them, which is
% what applying a gate to those indices looks like.
\newcommand{\tngate}[5][4mm]{%
  \tn@spanx{#4}%
  \pgfmathsetlength{\pgf@xb}{\tn@hi-\tn@lo+2*(#1)}%
  \pgfmathsetlength{\pgf@xc}{(\tn@lo+\tn@hi)/2}%
  \edef\tn@wd{\the\pgf@xb}\edef\tn@md{\the\pgf@xc}%
  \tn@lowy{#4}%
  \pgfmathsetlength{\pgf@yb}{\tn@lo-(#5)}%
  \node[tn gate, minimum width=\tn@wd] (#2) at (\tn@md,\the\pgf@yb) {#3};%
}
